\chapter{P0：先确认数据能不能放心用}

\section{承上启下}

在项目起点，我们还没有一批可以放心使用的数据，也没有资格讨论检索或 Agent。P0 的任务是回答一个更基础的业务问题：\textbf{投入下载和开发之前，来源、字段、本地资源与公开边界是否足以支持一个本地证据原型？}

\section{输入、边界与价值}

\textbf{输入}包括官方数据说明、评论与元数据的最低字段需求、预期文件规模、本地 Python 与磁盘条件，以及公开作品集的目标。用运营语言说，就是先确认“数据从哪里来、里面有没有后面要用的内容、我的电脑能不能处理、哪些内容不能公开”。

\textbf{边界}是只做有界预检，不下载完整数据，不检查评论内容质量，也不证明数据可以公开再分发。P0 带来的业务价值是提前暴露资源和合规风险；技术价值是把环境与来源假设变成可重复检查的证据。

\begin{boundarybox}
“公开可访问”只说明当前能够从公开端点取得响应，不等于数据附带再分发许可。任何评论衍生文本的公开展示仍需单独的数据使用审查；原始数据继续留在本地且不进入 Git。
\end{boundarybox}

\section{开发过程}

\subsection{1. 把模糊可行性拆成四个问题}

P0 分别检查来源可达性、字段可用性、运行环境和容量，并把官方事实、现场观察、设计决定、开放问题和阻塞项分开记录。这样，后来即使来源或环境变化，也能知道需要重查哪一项。这里的重点不是写复杂脚本，而是把“我觉得应该可以”变成几项能重新检查的事实。

\subsection{2. 先用很小的请求探测来源}

项目对评论和元数据来源各请求一个很小的字节范围。通俗地说，这像先确认仓库门能打开、标签与预期一致，而不是立即把整批货搬回家。观察到的部分内容响应与 gzip 类型，只支持“当时的端点可达且对象形态符合预期”这一有限结论。

\subsection{3. 检查本地环境和容量}

预检确认支持的 Python、操作系统和可用空间能够覆盖预计下载与后续处理需求。容量估算必须为压缩输入、临时文件和派生产物留出空间，不能只比较压缩包大小。

\subsection{4. 先写公开边界}

公开版本只允许代码、合成夹具、聚合结果和经审查的文档。原始评论、真实标识、密钥、大型索引和本地路径从一开始就被排除，而不是发布前临时清理。

\section{本阶段学到的知识}

\textbf{有界探测}用最小网络和存储成本验证关键假设；\textbf{容量预检}避免流程中途失败；\textbf{事实分层}防止把观察、决策和推断混在一起；\textbf{数据可用性与许可分离}则是公开项目的治理基础。换成初学者的话：先小范围检查，再决定是否值得投入时间和空间。

\begin{beginnerbox}
项目例子：两个来源都对一字节范围请求返回有效响应，只能说明 P0 当时观察到“可达”。它不能证明完整文件已经下载，不能证明内容正确，也不能代替许可证判断。
\end{beginnerbox}

\section{问题、修复与验证}

最危险的误解是“能下载，所以能公开”。修复不是猜测许可，而是缩小公开面，并在报告中保留“公开展示需另行审查”的开放问题。验证则重新执行有界来源探测、Python 与容量检查，确认没有当前阻塞项。

\textbf{小练习：}把“端点返回部分内容”改写成三种句子：可声称、不可声称、仍待核实。合格答案应把可达性、完整下载和再分发权分开。

\section{阶段产出与验收边界}

P0 的权威产出是 \filelink{docs/reports/phase-0-preflight.md}。其 Gate 证明的是：在记录时点，官方来源有界可达，环境和容量满足当前计划，公开边界已有设计。它不证明完整数据质量、再分发权、业务需求或任何商业效果。P0 在这一边界内为 \statuspassed。

\section{怎样进入下一阶段}

P1 只能使用 P0 核实过的两个来源角色，开始可恢复取得和完整性登记。若来源身份、所需字段或容量条件变化，应先重跑 P0，而不是直接把新文件塞进原始目录。

\section{面试怎么讲}

\begin{interviewbox}
“我没有直接下载数据或搭 Agent，而是先用有界请求验证两个官方来源，并检查 Python、容量和公开边界。这个阶段只证明项目在当前条件下可开展；公开可访问不等于有再分发许可，所以原始数据保持本地，评论衍生文本仍需单独审查。”
\end{interviewbox}

\textbf{掌握标准：}能够解释 P0 的四项检查，并明确说出至少三项它没有证明的结论。
